% Supplementary information for the city-square exposure paper.
% Build: pdflatex si && pdflatex si
\documentclass[journal]{IEEEtran}

\usepackage[T1]{fontenc}
\usepackage{amsmath,amssymb}
\usepackage{booktabs}
\usepackage[hidelinks]{hyperref}
\usepackage{cleveref}

\crefname{table}{Table}{Tables}

% Supplementary numbering, so nothing here collides with the paper.
\renewcommand{\thesection}{S\arabic{section}}
\renewcommand{\thetable}{S\arabic{table}}
\renewcommand{\theequation}{S\arabic{equation}}

% Same macros and the same notation as the paper.
\newcommand{\uloc}{\hat{u}}
\newcommand{\uext}{\hat{v}}
\newcommand{\nhat}{\hat{n}}
\newcommand{\khat}{\hat{k}}
\newcommand{\x}{\mathbf{x}}
\newcommand{\dOmega}{\,\mathrm{d}\Omega}

\begin{document}

\title{Supplementary information for ``City-square geometry sets the
distribution of relative pedestrian exposure at 15\,GHz''}
\author{Robin~Wydaeghe%
\thanks{The author is with the Department of Information Technology, Ghent
University / imec, 9052 Ghent, Belgium.}}

\maketitle

This document holds the two procedures the paper defers to it, the registration
of a panorama against the mesh and the projection of image labels onto
triangles, together with the seeds and ray counts behind every reported number.
Three further sections give the per-square numbers behind the headline figures,
the material assignment, and the convergence sweeps at the resolution the body
had no room for. Notation, symbols and units follow the paper. Sections, tables
and equations here carry an S prefix, so a reference without one points at the
paper.

\section{Panorama registration}
\label{sec:si-registration}

A panorama arrives with a horizontal position and a heading from its provider
and with no altitude. Registration refines that pose by matching the sky
boundary the segmented image sees against the angular upper envelope of the
photogrammetric mesh. The skyline is used rather than the whole image because
tile noise at ground level is then absent from the objective. This section gives
the objective, the free parameters and their bounds, the seed study that
supplies the pose uncertainty, what the reported residual does and does not
measure, and which stations are usable.

\subsection{The objective}

The observed skyline is read from the segmented panorama. For each of 1024
equally spaced image columns, the lowest row belonging to the sky class is taken
over the top 72\% of the rows. A column with no sky at all has no skyline and is
dropped rather than filled, so no observation is fabricated. A median filter of
width 11 columns then removes the gaps a cable or a bare tree punches in
otherwise connected sky. When the segmentation manifest supplies a list of
structural classes, a column is kept only when the pixel just below its boundary
carries one of them, which are building, wall, bridge and tunnel. Each kept
column becomes a unit ray in the panorama frame.

The modelled skyline is read from the mesh vertices. Vertices more than 8\,m
away in horizontal range are sorted into the same 1024 azimuth bins, and each
bin takes the largest elevation among its vertices. Bins with no vertex are
filled by interpolation across the wrap. A percentile filter of width 11 bins
then takes the 90th percentile of each neighbourhood.

That percentile is the one setting in the objective with a physical
consequence. At the 50th percentile the filter is an azimuthal median, which
rejects the isolated high vertices photogrammetry leaves behind and also shaves
every convex roof peak, and a roofline is convex upward almost everywhere it
matters. The modelled skyline then sits too low, and the optimizer has one cheap
way to raise every modelled rooftop at once, which is to lower the camera.
Measured at the camera's own altitude, the observed skyline sat $3.83^\circ$
above the modelled one at Korenmarkt and $1.56^\circ$ at Milan under the median
filter, falling to $1.11^\circ$ and $0.52^\circ$ under the 90th percentile.

The cost is the mean absolute angular difference between the observed and the
modelled elevation over the kept samples, made robust by clipping every residual
above the 0.8 quantile to that quantile, which is what keeps trees and cranes
from steering the fit. A small prior on the orientation offsets is added,
$0.001\,[(\Delta\theta/3)^2 + (\Delta\phi/3)^2]$ with both offsets in degrees.
The residual quoted throughout is that robust mean converted to degrees.

\subsection{Free parameters, bounds and optimizer}

Six parameters are solved for: two horizontal offsets, an altitude offset, a
yaw correction, and pitch and roll corrections about the provider's values.
Their bounds are $\pm4$\,m horizontally, $-3.0$ to $+3.0$\,m in altitude,
$\pm20^\circ$ in yaw and $\pm6^\circ$ in pitch and roll.

A seventh parameter exists and is switched off in every run reported here. It
adds a constant elevation bias to the modelled skyline, and it is degenerate
with altitude: a roofline too low by $d$ meters and a camera too high by $d$
meters produce the same angular error $d\cos^2\!\alpha / R$ at every azimuth.
Fitting both at Korenmarkt traces a straight valley at $1.47^\circ$ of bias per
meter of altitude, along which three meters of altitude cost $0.05^\circ$ of
residual. The bias is therefore kept as a diagnostic and the altitude is
measured against the support mesh instead of being inferred.

The fit is a global one by differential evolution~\cite{storn}, with a
population size factor of 7, at most 600 generations, immediate population
updating, an absolute convergence tolerance of $10^{-7}$ on the cost, and a
local polish at the end~\cite{scipy}. Before the fit runs, the mesh vertices are
pruned to those that can reach the skyline anywhere the optimizer is allowed to
go. The prune is the union over the corners and the center of the translation
search box of the vertices standing within $4^\circ$ of the envelope seen from
that point, with every fortieth vertex retained unconditionally. Pruning at the
initial camera alone would be wrong, because a vertex hidden from the start
position can be the silhouette from a meter away.

\subsection{The seed study, and what the residual is not}

Differential evolution is stochastic, so the fit is repeated from eight
independent seeds on identical data. The seeds are 714261, 1, 2, 3, 5, 8, 13 and
21. The lowest-cost fit of the eight is the pose that ships. The sample
covariance of the fitted parameters across the eight is what ships beside it as
the pose uncertainty, and the horizontal one-sigma figure quoted per site in
\cref{tab:si-registration} is the square root of the sum of the two horizontal
variances of that covariance.

The residual is a goodness of fit between two silhouettes and nothing else. It
is not an error bar on the pose, which is why every pose carries the seed
ensemble as well. Three failures pass it, and each is caught by a different
test. A camera driven inside the geometry scores a low residual, because the
silhouette it is matching is the inside of a wall and it matches it well. A site
whose skyline is uniformly high carries almost no orientation information, so a
low residual there says little about the heading. A pose can also be several
meters from where the camera stood while still matching the horizon, which
leaves the residual unchanged and the surfaces the camera is pointed at wrong.

\subsection{Admission, and the second test}

A pose is admitted when its residual is at most $4.0^\circ$, its sky conflict is
at most 0.5, and the median range of any conflict is at least 2.0\,m.

The sky conflict is the second test and it shares no code with the objective. A
$512 \times 256$ equirectangular grid of directions is cast from the fitted
camera against the support mesh, dropping elevations below $-60^\circ$ because
the rig and the road directly under it are not a registration signal. The
conflict is the fraction of directions the segmentation calls sky for which the
mesh returns a first hit, and it is read together with the median range of those
hits. A healthy pose sits near zero with its few conflicts tens of meters away.
A pose at 1.0 with a median conflict range under a meter is inside a building.

As the poses were first shipped, 26 of the 83 were inside the geometry by that
test and 6 of those 26 passed the residual gate, so the two tests are not
redundant and the residual on its own admits poses that cannot be used.

The mechanism behind the 26 is the crop, and it is measured rather than
inferred. Every registration in the study was fitted against the 130\,m mesh,
which the candidate vertex count stored in each pose reproduces exactly at all
six registered sites and at no other radius: 23\,640 at Grand-Place, 38\,138 at
Plaza Mayor, 36\,782 at the Zocalo, 47\,184 at Old Town Square, 38\,304 at Times
Square and 27\,371 at Hachiko. A 130\,m crop is the whole skyline of a low-rise
square and a fraction of the skyline of a high-rise canyon. Truncating it lowers
the modelled skyline, the optimizer buys the residual back by sinking the
camera, and it keeps doing so until it reaches the bound. All 25 of the poses
called inside the geometry that carry a recorded dive dropped more than 1.5\,m,
with a median dive of 2.97\,m against a 3.0\,m bound, whereas the 41 admitted
poses have a median dive of 0.58\,m and only one in eight goes past 1.5\,m.

Scoring the untouched initial pose against both crops separates the sites
cleanly. Taking the median over each site's first three stations, the extra crop
is worth $-0.00^\circ$ at Old Town Square, $+0.00^\circ$ at Plaza Mayor,
$-0.02^\circ$ at Grand-Place, $+0.48^\circ$ at the Zocalo, $+1.37^\circ$ at
Times Square and $+5.60^\circ$ at Hachiko. The shipped crop was therefore
adequate at four sites, marginal at a fifth and wrong at one.

Hachiko was re-registered against the 250\,m mesh with the altitude bound
clamped to $-1.5$ to $+3.0$\,m, which keeps the optical center at least a meter
above the ground measured under that camera and excludes no pose that was ever
any good. Its three stations move from $6.53^\circ$, $4.80^\circ$ and
$6.19^\circ$ to $1.42^\circ$, $3.70^\circ$ and $3.08^\circ$, and their sky
conflicts from 1.00 at 0.5, 0.3 and 0.1\,m to 0.02, 0.07 and 0.01 at 108.1, 20.6
and 33.9\,m. Zero usable stations become three. The two changes do different
work: the wider crop moves the residual, from a median of $6.19^\circ$ to
$3.08^\circ$, and the clamp moves the sky conflict at a cost of $0.34^\circ$ of
residual at one station and nothing at the other. Across the study the
re-registrations take the count of poses inside the geometry from 26 to 14 and
the admitted count to 51, and 6 of the 14 still pass the residual gate, so the
disagreement between the two tests survives the fix.

Times Square keeps a residual near $9^\circ$ on the wider crop and the crop is
not its problem. Scanning yaw over the entire circle at the initial pose moves
the residual by under $2^\circ$ at every one of its 14 stations, and four of
them score their minimum more than $120^\circ$ away from the provider's heading,
one at $178^\circ$. The control says this is the site and not the method: the
same scan puts 13 of 14 stations within $2^\circ$ of the provider heading at Old
Town Square and all three within $2^\circ$ at Hachiko. The modelled skyline at
Times Square sits at a median elevation of 55 to $62^\circ$ in every direction.
A uniformly high skyline has no features, and rotating a camera inside one
changes almost nothing. Skyline registration needs a ragged horizon, and the
deepest canyon in the set is the one square without one. The yaw scan behind
this paragraph is recorded in the repository notes and no output file was kept
for it, so it can be repeated but not read off disk.

\subsection{First-interaction coverage is a sharper test than the residual}
\label{sec:si-coverage-test}

A third test came from the propagation side and it is the one the paper quotes.
Trace the reference square at each admitted station position and record the
share of the power incident at the first interaction that lands on a triangle
some panorama photographed. The run is Korenmarkt at the 130\,m crop, 15\,GHz,
200\,000 rays per standpoint. \Cref{tab:si-station-coverage} gives the eight
admitted stations of that binding, ordered by their skyline residual.

Seven of the eight sit between 0.978 and 0.999 at the first interaction, which
confirms the geometric argument of the paper to about two parts in a thousand.
The eighth returns 0.362. It carries the largest residual in the admitted set at
$3.67^\circ$, but the residual does not separate it from the rest, because the
admitted set runs from $1.08^\circ$ to $3.08^\circ$ against a gate at
$4.0^\circ$ and the station at $3.06^\circ$ scores 0.998. The low station is not
somewhere nobody walked either. It sits 9.1\,m from a station scoring 0.999, and
8.1\,m from it once both are taken at their registered rather than their
provider positions.

One detail of how the probe was run has to be given, because it changes what the
row means. The eight probe standpoints of this run are the positions the
provider supplied, not the positions registration returned, while the set of
photographed triangles they are scored against was built from the registered
poses. The eight cameras moved 3.91 to 5.50\,m horizontally and $-1.61$ to
$+1.37$\,m vertically during registration, so every one of the eight probes is a
few metres away from the camera whose photographs it is being scored against.
The size of that gap does not order the result. The camera that moved furthest,
5.50\,m, scores 0.978, and the one that scores 0.362 moved 5.28\,m.

The same camera reads 1.000, 0.962 and 0.919 over the first three interactions
in the 250\,m binding, in line with every other station. That comparison is
suggestive and it is not controlled. Three things differ between the two runs at
once: the probe standpoints are the registered positions rather than the
provider ones, the crop radius is 250\,m rather than 130\,m, and the station set
holds nine cameras rather than eight. No single one of them can be credited with
the change on the evidence here.

What the two rows do support is narrow and is what the paper claims. At one
fixed binding, one standpoint out of eight returned 0.362 where the rest
returned 0.978 to 0.999, and the skyline residual gave no warning. A
first-interaction coverage probe therefore finds something the residual does
not, and it costs one trace. Why that standpoint is the one is not resolved
here. The station is excluded from the coverage figures the paper quotes.

The residual 2\% at the seven good stations is the resolution of the ray grid
rather than the pose. A triangle small enough to fall between adjacent rows of a
1536-row equirectangular cast collects no ray and is therefore not observed,
whatever the camera saw.

\begin{table}[!t]
  \caption{First-interaction coverage at the eight admitted Korenmarkt stations,
  130\,m crop, ordered by skyline residual}
  \label{tab:si-station-coverage}
  \centering
  \begin{tabular}{lcccc}
    \toprule
    Station & Residual & \multicolumn{3}{c}{Share of incident power on} \\
     & (deg) & \multicolumn{3}{c}{photographed triangles} \\
    \cmidrule(l){3-5}
     & & first & second & third \\
    \midrule
    1084407470281938 & 1.08 & 0.978 & 0.909 & 0.823 \\
    582445657694750  & 1.30 & 0.998 & 0.952 & 0.890 \\
    1019442960256615 & 2.15 & 0.997 & 0.962 & 0.897 \\
    3367014310197013 & 2.32 & 0.995 & 0.959 & 0.901 \\
    706535575184668  & 2.39 & 0.996 & 0.958 & 0.902 \\
    1419513849204492 & 3.06 & 0.998 & 0.965 & 0.915 \\
    986493819697753  & 3.08 & 0.999 & 0.963 & 0.919 \\
    1304023184185511 & 3.67 & 0.362 & 0.351 & 0.226 \\
    \addlinespace
    1304023184185511, & 3.67 & 1.000 & 0.962 & 0.919 \\
    \quad at the 250\,m binding & & & & \\
    \bottomrule
  \end{tabular}
\end{table}

\subsection{Which stations are usable}

\Cref{tab:si-registration} gives the whole corpus. Eighty-three panoramas were
acquired across the eleven squares and all 83 registered. Fifty-one poses are
admitted.

Three squares carry no panoramas at all. At Krakow and Toulouse the capture with
the most stations is an indoor virtual tour, 162 panoramas of the Cloth Hall
interior and 118 of the interior of the Capitole, and a downward cast at each
candidate finds open sky above 0 of 162 and 0 of 118 of them. London's screened
capture is outdoors, with 111 of 142 stations under open sky, and no panoramas
were acquired there.

One square registered and lost everything at admission. All 14 Times Square
poses registered and none is admitted, for the reason given above.

Two squares carry the material axis. Twelve panoramas in the corpus carry a
material label rather than only an entity label, eleven of them at Korenmarkt
and one at Milan. Every panorama at the other nine squares carries the entity
axis only. This is why the material comparison in the paper is a one-square
result and why the eleven-square comparison runs on geometric classes.

\begin{table*}[!t]
  \caption{Registration, per square. Residuals are the robust skyline residual
  in degrees. Pose sigma is the horizontal one-sigma of the eight-seed
  ensemble covariance. Admission requires a residual of at most $4.0^\circ$, a
  sky conflict of at most 0.5, and a conflict range of at least 2.0\,m}
  \label{tab:si-registration}
  \centering
  \begin{tabular}{lcccccccc}
    \toprule
    Square & Panoramas & Registered & Median & Worst & Median pose & At altitude & Inside the & Admitted \\
     & & & residual & residual & sigma (m) & bound & geometry & \\
    \midrule
    Brussels, Grand-Place        & 14 & 14 & 2.87 & 11.16 & 0.47 & 6 & 6 & 8 \\
    Ghent, Korenmarkt            & 13 & 13 & 3.06 & 8.10  & 0.18 & 7 & 0 & 9 \\
    Krakow, Rynek Glowny         & 0  & 0  & --   & --    & --   & 0 & 0 & 0 \\
    London, Trafalgar Square     & 0  & 0  & --   & --    & --   & 0 & 0 & 0 \\
    Madrid, Plaza Mayor          & 10 & 10 & 0.40 & 2.92  & 0.26 & 5 & 4 & 6 \\
    Mexico City, Zocalo          & 14 & 14 & 2.76 & 4.60  & 1.86 & 4 & 0 & 12 \\
    Milan, Piazza del Duomo      & 1  & 1  & 0.75 & 0.75  & 0.26 & 1 & 0 & 1 \\
    New York, Times Square       & 14 & 14 & 9.13 & 11.25 & 1.44 & 5 & 3 & 0 \\
    Prague, Old Town Square      & 14 & 14 & 0.82 & 9.33  & 0.25 & 2 & 1 & 12 \\
    Tokyo, Hachiko crossing      & 3  & 3  & 3.08 & 3.70  & 0.48 & 1 & 0 & 3 \\
    Toulouse, Place du Capitole  & 0  & 0  & --   & --    & --   & 0 & 0 & 0 \\
    \midrule
    Total                        & 83 & 83 &      &       &      & 31 & 14 & 51 \\
    \bottomrule
  \end{tabular}
\end{table*}

\section{Projection onto triangles}
\label{sec:si-projection}

Two routes turn labelled pixels into per-triangle evidence. One casts a ray grid
outward from the camera and tallies a class per triangle of the mesh the tracer
will use. The other cuts each visible triangle in the image against the label
boundaries and carries the pieces back onto the triangle's own plane. Both are
described here, because the exposure runs read the first and the coverage
figures quote both.

\subsection{The station cast}

From each admitted station an equirectangular ray grid of 1536 rows is cast
against the mesh of the run. Per triangle the cast records the modal semantic
class it collects and the number of rays that survived transient rejection, with
one row per station, so a triangle is observed when at least one station landed
a ray on it.

The grid resolution sets a floor on what can be covered, because a triangle no
ray landed on is a triangle no station saw. At Plaza Mayor over the 130\,m crop
the same six admitted stations bind 18.04\% of the triangle area at 384 rows and
19.43\% at 768, so a fourfold cheaper grid costs about 7\% of the answer in
relative terms. Every per-site binding records a grid height of 1536 rows, and
the grid is held fixed across squares because the covered fraction is compared
between squares.

One binding is not on that grid and it is the one the image-derived exposure
numbers rest on. The Korenmarkt 130\,m binding that carries the material axis is
a separate file, and its sidecar records no grid height at all. It was built at
1024 rows, which was established by rebuilding at 1024 and reproducing the
stored ray, transient, clean-ray and modal-class arrays bit for bit. A rebuild
at 1536 casts 2.25 times as many rays and moves the modal class on 4.9\% of the
triangles both grids see. What that costs at the end of the chain was measured
on the nine-station variant, where a uniform 1024 rebuild agrees with the stored
file on 99.45\% of the 2559 triangles at issue, moves the bound area from
11.019\% to 10.982\% and moves the evidence-ladder shift from 0.350 to
0.352\,dB isotropic and from 0.364 to 0.368\,dB rooftop. The grid height of the
shipped file is not recorded in it, and that is a provenance defect rather than
a numerical one.

The class the cast collects is a semantic entity rather than a material. It is
mapped onto the radio material vocabulary by a table that is a property of the
vocabulary and not a site measurement. All 83 station-level segmentation files
carry the same 65-class vocabulary from the same segmentation
checkpoint~\cite{mask2former,vistas}, which is checked label by label rather
than assumed, so the same table is passed through to every square.

\subsection{The cutter}

The cutter reverses the direction of travel. Rather than tiling the label image
and splatting tiles onto the mesh, it projects each support triangle into the
image, cuts it there against the label boundaries, and maps the pieces back onto
the triangle's supporting plane.

Within one rectilinear crop the camera is an ordinary pinhole, so a triangle
maps to a triangle under a projective map that is invertible on the triangle's
interior. The inverse used is the exact intersection of the pixel ray with the
supporting plane, not barycentric interpolation of the projected corners, which
is wrong under perspective. Checked against an independent first-hit depth
buffer at Korenmarkt, unprojecting all 707\,620 first-hit pixels of one crop and
testing the result against the plane of the face the buffer reports gives a
median residual of $4.7 \times 10^{-7}$\,m, a 95th percentile of
$3.7 \times 10^{-6}$\,m and a worst case of $1.9 \times 10^{-5}$\,m.

Two guards are not optional. Each triangle is clipped in camera space against a
near plane at 0.05\,m before it is projected, because at Korenmarkt yaw 0 four of
the 1431 visible support triangles have a vertex behind the camera plane and
those four own 33.5\% of the hit pixels in the view. And the projection is a
bijection on the visible surface only, since two triangles routinely project to
the same pixel, so a first-hit triangle identifier buffer decides ownership and
any piece a different triangle owns is rejected rather than painted.

The cut runs in four steps. First, every pixel receives one reason code:
paintable, no mesh hit, low confidence, transient object, clutter in front, mesh
or pose conflict, or not a support surface. Distance and class are kept as
separate causes on purpose, because a person's depth often agrees with the wall
they stand against and only the class can reject them. Second, paintable pixels
are grouped into connected components that share a class and are continuous in
first-hit range, with a relative step test so grazing surfaces stay whole.
Folding depth continuity into the grouping is what puts occlusion silhouettes
into the output, since a wall seen past a nearer building becomes a different
component from that building even when both carry the same label. Components
below 64 pixels merge into the depth-continuous neighbor with the longest shared
boundary. Third, boundaries are traced on the pixel-corner crack grid rather
than per class mask, so the boundary between two components is one shared
object, and each chain between junctions is simplified once by the
Douglas-Peucker rule~\cite{douglas} with its junctions held fixed, which can
open neither a gap nor an overlap. Fourth, each visible triangle is clipped to
the near plane and the crop rectangle and rasterized. When its footprint lies in
one component and the triangle owns at least 98\% of it, the triangle is emitted
untouched, because the support triangulation is meter-scale photogrammetry and
subdividing it where the labels do not change buys nothing. Otherwise the local
chains are noded against the triangle outline, the resulting arrangement faces
are depth-tested and support-tested, and each survivor is triangulated by ear
clipping.

\subsection{Crop geometry}

Each panorama is cut through four rectilinear crops at yaw 0, 90, 180 and
$270^\circ$ with zero pitch, each $90^\circ$ across and $1024 \times 1024$
pixels. Pixel row $i$ and column $j$ has its center at image coordinate
$(j + 0.5, i + 0.5)$.

The boundary simplification tolerance is the one crop-side parameter that was
swept. \Cref{tab:si-tolerance} gives that sweep at Korenmarkt yaw 0. The face
count falls steeply up to about 1.5 pixels and then flattens, because past that
point the output is dominated by the support triangulation rather than by the
boundary description. The shipped tolerance of 1.5 pixels is $0.168^\circ$ of
worst-case boundary displacement at the crop center, where a $90^\circ$ crop has
its coarsest angular pixel pitch. That is 29\,mm at 10\,m of range, about 1.5
wavelengths at 15\,GHz, which is fine enough for material assignment and not
fine enough for coherent phase on a path that hugs a boundary. Raising the
minimum component size from 64 to 256 pixels nearly doubles the class change
while saving 2\% of the faces, so 64 is the default.

\begin{table}[!t]
  \caption{Boundary simplification sweep at Korenmarkt yaw 0. Coverage is the
  share of paintable pixels an output face lands on after a round trip through
  the image, and class change is the share of those whose label differs}
  \label{tab:si-tolerance}
  \centering
  \begin{tabular}{cccccc}
    \toprule
    Tolerance & Minimum & Faces & Angular & Class & Coverage \\
    (px) & island (px) & & (deg) & change & \\
    \midrule
    0.0 & 64  & 12\,191 & 0.000 & 0.077\% & 99.23\% \\
    1.0 & 64  & 3\,987  & 0.112 & 0.150\% & 99.02\% \\
    1.5 & 64  & 3\,306  & 0.168 & 0.190\% & 98.91\% \\
    2.5 & 64  & 2\,932  & 0.280 & 0.277\% & 98.72\% \\
    4.0 & 64  & 2\,768  & 0.448 & 0.353\% & 98.55\% \\
    1.5 & 256 & 3\,252  & 0.168 & 0.348\% & 98.91\% \\
    \bottomrule
  \end{tabular}
\end{table}

The four Korenmarkt crops give 10\,534 faces against 257\,310 from the
tile-splatting path they replace on the same mesh, pose and label maps, a factor
of 24. The floor is hard rather than tunable: a triangle inside one component is
preserved, so the output can never fall below the number of visible support
triangles, which is 1431 at yaw 0. At 3306 the cutter sits 2.3 times above that
floor, and going lower means decimating the support mesh.

\subsection{The per-face record}

Each face carries an indexed triangle in world meters together with the source
triangle it was cut from, its outward normal, which points back toward the
capture point by construction, its centroid, its metric area, its exact solid
angle from the capture point, its class, the class posterior over the taxonomy,
a confidence, a pixel support count, a visible fraction, and a compressed index
of the source pixels it was built from. A material posterior is written by the
same code when a material model was run, and it is absent from the shipped
Korenmarkt fishnets, which carry the entity taxonomy only. Candidates the cut
refused are kept in a parallel table with a reason code and their image area, so
what was removed and why remains readable downstream. Image area is conserved
exactly: at every view the emitted and rejected areas sum to the clipped source
area to the last pixel.

\subsection{Visible fraction and confidence}

The visible fraction of a face is the share of its rasterized footprint that the
first-hit identifier buffer assigns to that face's own source triangle. A piece
owning less than 0.5 of its footprint is rejected rather than kept at a low
visible fraction, so the kept faces are a truncated sample and that threshold is
reported beside the occlusion budget for exactly that reason.

Confidence is built from the pixels a face owns. The class is the
confidence-weighted majority over those pixels, the class posterior is that
weighted histogram normalized, and the confidence is the winner's share of the
histogram multiplied by the mean per-pixel confidence of the winning pixels.
Where a material model is present the material posterior is the mean of the
per-pixel material probabilities over the same winning pixels.

The station cast keeps no such record. Its per-triangle output is the modal
class, the number of rays, the number of rays lost to transient objects and the
number that survived, one row per station. There is no per-triangle confidence,
no visible fraction, no incidence angle and no record of which station won. The
per-pixel confidence raster produced by the segmentation network is present in
every station file and the cast does not read it, so no confidence gate acts
anywhere on this route. Its threshold for observing a triangle is one
transient-free ray. When several stations see the same triangle the winner is
chosen by a vote weighted by each station's transient-free ray count on that
triangle, which stands in for view quality because a triangle seen edge on
collects few rays.

The occlusion budget reports what occlusion removed as a share of the projected
image area of the visible support triangles, split three ways because the three
fail differently. The within-cell term is area an accepted face carries but does
not own, the absent term is whole cells rejected with nothing downstream to
replace them, and the deferred term is whole cells handed to the separate
pedestrian layer. Over the four Korenmarkt crops on the entity taxonomy, at
10\,534 faces, those are 0.725\%, 9.169\% and 20.490\%, totalling 30.385\%. On
the 17-class material taxonomy, at 16\,451 faces, they are 0.474\%, 9.663\% and
27.252\%, totalling 37.389\%.

Fidelity is measured by rasterizing the output faces back into the image with a
nearest-surface test and comparing against the source label map.
\Cref{tab:si-roundtrip} gives that round trip per view. Both directions of the
visibility error are reported because they fail differently for propagation: a
phantom face invents a scatterer where the camera cannot see one, and a deleted
face removes a path.

\begin{table}[!t]
  \caption{Round trip of the cut surfaces back into the image, per crop at
  Korenmarkt. Leaked pixels are unpaintable pixels an output face nonetheless
  paints, which is the error of painting a person onto the wall behind them}
  \label{tab:si-roundtrip}
  \centering
  \begin{tabular}{lccccc}
    \toprule
    View & Class change & Coverage & Leaked & Phantom & Deleted \\
    \midrule
    yaw 0   & 0.190\% & 98.91\% & 0.294\% & 0.380\% & 1.086\% \\
    yaw 90  & 0.126\% & 98.74\% & 0.117\% & 0.094\% & 1.262\% \\
    yaw 180 & 0.123\% & 98.92\% & 0.121\% & 0.129\% & 1.080\% \\
    yaw 270 & 0.250\% & 99.27\% & 1.111\% & 0.063\% & 0.731\% \\
    \bottomrule
  \end{tabular}
\end{table}

Withheld transient pixels are not painted by anything. They leave a hole in the
static surface, and the hole is filled by the separate pedestrian layer rather
than by the cutter, which is why the occlusion budget keeps the deferred term
apart from the absent one. On the four Korenmarkt crops that layer holds 18
people, of stature 1.40 to 1.71\,m with a mean of 1.56\,m, standing 3.2 to
15.0\,m from the capture point.

\subsection{How much of a square the projection binds}

The two routes are built from the same photographs by different machinery and
they agree to a few points of area without being the same number. Over the
130\,m crop the station cast binds 31.1\%, 26.2\%, 20.2\% and 13.3\% of the
triangle area at Old Town Square, the Zocalo, Plaza Mayor and Grand-Place,
against 33.0\%, 23.4\%, 18.7\% and 11.2\% for the cutter, so the cutter is the
larger of the two at one square and the smaller at the other three. Which is
nearer the truth is not settled here and should not be read off the fractions.

At the 250\,m crop the headline runs use, the station cast binds 9.9\% at Old
Town Square, 8.1\% at Hachiko, 7.0\% at the Zocalo, 5.1\% at Plaza Mayor, 4.6\%
at Grand-Place, 4.0\% at the Duomo and 3.2\% at Korenmarkt. The spread is not a
quality signal. Bound area rises with the number of admitted stations and falls
with the size of the crop relative to the square, so the Duomo reaches 4.0\%
from a single camera because the facade it faces is large and close, and Hachiko
reaches 8.1\% from three stations for the same reason in reverse. The admitted
cameras at every square with a station set sit within 29 to 62\,m of their own
centroid, against a crop that reaches 250\,m, so most of a crop is roofs and rear
elevations no street-level capture in this study has seen.

\section{Seeds and ray counts}
\label{sec:si-seeds}

Every susceptibility in the paper is a Monte Carlo estimate. This section gives
what a seed controls, how seeds are assigned, how many rays each result was
given, how the error on a reported number was measured, and which comparisons
are paired and therefore carry less noise than their two sides do.

\subsection{What the seed fixes}

One standpoint is traced by one call, and that call opens a single NumPy
generator, a PCG64 instance initialised from an integer seed through
SeedSequence \cite{numpy}. That one generator supplies every random number the
trace uses. There are three of them. The launch directions are drawn uniformly
on the sphere. At each surface interaction the choice between the specular and
the diffuse branch is a draw against the Rayleigh specular share. A ray taking
the diffuse branch is given a cosine-weighted direction about the local normal.
Russian roulette would be a fourth, but at the operating point of the paper it
never fires, because \texttt{roulette\_start} is set one past the interaction
budget.

Nothing else in a trace is random. Geometry, material assignment, the ground
datum and the standpoint list are all fixed before the generator is opened, so
two traces with the same seed and the same scene return bit-identical numbers.
That was checked rather than assumed: in the seed study of
\cref{sec:si-seeds-error} the replica that reuses the published stream
reproduces the stored file to zero difference at all 120 standpoints and all
three illumination models.

The launch direction of an adjoint ray is the local arrival direction
$\uloc$, so it is also the axis $\rho$ is binned on. Each drawn direction is
assigned to the nearest of the $M = 512$ near-equal-solid-angle cells, and the
cell value is the mean over the rays assigned to that cell rather than a sum
divided by $N/M$. A cell that received no ray contributes nothing instead of
dividing by zero. This is a stratification in $\uloc$ and it does not reduce the
variance that matters, because the weight $Q$ is evaluated on the escape
direction and the escape direction is not the launch direction.

\subsection{How seeds are assigned}

The eleven-square runs and every ablation built on the same walk use one base
seed of 7 and give standpoint $i$ the seed $7 + 1000\,i$, where $i$ is the index
of the standpoint in the site's walk rather than its row in the output file.
Two consequences follow and both are used. No two standpoints inside one run
share a stream, so the standpoints of a run are independent draws. Two runs over
the same walk use the same stream at each matching standpoint, so an ablation
that changes only the material assignment or only the interaction budget is a
paired comparison at every standpoint.

The smaller studies were written as separate scripts and do not all follow that
map. The law comparison and the sub-street ablation use base seed 0 with an
offset of one per standpoint. The monostatic cross-check uses base seed 7 with
an offset of 1000 for the walk sweep and 977 for its station sweep. The material
posterior ablation uses base seed 0 with an offset of 1000, and draws its
material realisations from a separate generator seeded at $1000 + j$ for
realisation $j$. The bystander noise floor uses $7 + 1000\,i + 131\,r$ for
replica $r$.

The crop convergence sweep is the one exception worth flagging. It passes no
per-standpoint seed, so every one of its 32 observers at every one of its nine
radii is traced with the configured seed of 11. The 32 observers within one
radius are therefore not independent of one another, and the same launch set is
reused at all nine radii. That makes the sweep a paired comparison across
radius, which is what it is for, and it means the scatter between its rows is
not a Monte Carlo error estimate.

Distinct integer seeds passed through SeedSequence are not covered by a formal
non-overlap guarantee. The generator state space makes an overlap between two
streams of a few million draws a negligible risk, and no stronger claim than
that is made here.

\subsection{Ray counts}

\cref{tab:si-rays} lists what each result was given. Rays are per standpoint.
The interaction budget is the maximum number of surface interactions a ray may
make before it is dropped.

\begin{table*}[t]
  \centering
  \caption{Ray count, interaction budget and standpoint count behind each
  result. Rays are per standpoint. Two entries are deterministic and shoot no
  random rays.}
  \label{tab:si-rays}
  \footnotesize
  \begin{tabular}{llrlrrl}
    \toprule
    Result & Run tag & Sites & Standpoints & Rays & $L$ & Crop \\
    \midrule
    Eleven squares, headline & \texttt{city250\_L3} & 11 & 80 & 200\,000 & 3 & 250\,m \\
    Eleven squares, four interactions & \texttt{city250\_corrected} & 11 & 80 & 200\,000 & 4 & 250\,m \\
    Ground datum control & \texttt{city250\_datum} & 9 & 80 & 200\,000 & 4 & 250\,m \\
    Image evidence ladder & \texttt{*\_L3} & 1 & 120 & 200\,000 & 3 & 130\,m \\
    Crop convergence & \texttt{crop\_convergence} & 1 & 32 & 200\,000 & 3 & 9 radii, 60 to 340\,m \\
    Interaction budget sweep & \texttt{bounce\_evidence} & 4 & 40 & 200\,000 & 1 to 8 & 130 and 250\,m \\
    Seed spread and roulette cost & \texttt{bounce\_evidence} & 1 & 4, 8 seeds each & 200\,000 & 3 & 250\,m \\
    Illumination law, one square & \texttt{law\_comparison} & 1 & 40 & 150\,000 & 4 & 130 and 250\,m \\
    Sensitivity harvest & \texttt{harvest\_*} & 11 & 40 & 200\,000 & 4 & 250\,m \\
    Illumination law, eleven squares & \texttt{law\_ordering} & 11 & 40 & reweighted harvest & 4 & 250\,m \\
    Sub-street ablation & \texttt{*\_substreet} & 2 & 40 & 120\,000 & 3 & 250\,m \\
    Beamforming aperture sweep & \texttt{antenna11} & 11 & 32 & 150\,000 & 3 & 250\,m \\
    Beamforming budget control & \texttt{antennab4} & 1 & 32 & 150\,000 & 4 & 250\,m \\
    Beamforming steering artefact & \texttt{paper\_artefact} & 1 & 20 & 30\,000 & 3 & 250\,m \\
    Bystanders & \texttt{bystander\_study} & 1 & 12 & 120\,000 & 3 & 250\,m \\
    Bystander absorber control & \texttt{*\_absorber} & 1 & 8 & 120\,000 & 3 & 250\,m \\
    Bystander noise floor & \texttt{*\_noise\_floor} & 1 & 12, 3 seeds each & 120\,000 & 3 & 250\,m \\
    Monostatic return & \texttt{mono250\_*} & 11 & 40 & 100\,000 & 3 & 250\,m \\
    Material posterior ablation & \texttt{material\_vlm} & 1 & 24 & 60\,000 & 12 & 90\,m walk \\
    External cross-validation & \texttt{cross\_validation} & 2 & 8 & 200\,000 and 150\,000 & 4 & 250\,m \\
    Diffraction bound & \texttt{diffraction\_bound} & 3 & 20 & deterministic & -- & $720 \times 600$ grid \\
    \bottomrule
  \end{tabular}
\end{table*}

The headline run is 11 sites at 80 standpoints, so 880 traces of 200\,000 rays
each, which is $1.76 \times 10^{8}$ launched rays in total. Per-site wall clock
ran from 177\,s at Old Town Square to 730\,s at Grand-Place, and the whole
eleven-square run took 1.20 hours on one machine. The spread between sites is
scene size and not standpoint count, which is 80 everywhere.

Two rows of \cref{tab:si-rays} need reading with care. The eleven-square law
ordering shoots no rays of its own. It reweights the escape records already
stored by the sensitivity harvest, so both laws are evaluated on the same rays
at every standpoint and their difference carries no Monte Carlo noise at all.
The diffraction bound is not a Monte Carlo estimate either. It sweeps a fixed
grid of 720 azimuths by 600 elevations, logarithmic in elevation from a floor of
0.05 degrees, and so has no seed and no sampling error.

\subsection{How the error was measured}
\label{sec:si-seeds-error}

The error on a reported susceptibility is measured by retracing the same
standpoints under seed streams that share nothing, not by any within-run
variance estimate. Three such measurements exist and they are consistent with
one another.

The largest is the study behind the numbers the paper quotes. The 120 published
Korenmarkt standpoints of the geometric rung were retraced over eight seed
streams at 200\,000 rays, 512 cells, six surface interactions and the 130\,m
crop. \cref{tab:si-noise} gives the resulting standard deviations. The first
three columns are per standpoint and the last is the standard deviation of the
median over the 120 standpoints, which is the unit the paper actually reports.

\begin{table}[t]
  \centering
  \caption{Monte Carlo standard deviation from eight seed streams over 120
  Korenmarkt standpoints at 200\,000 rays. All entries in dB.}
  \label{tab:si-noise}
  \begin{tabular}{lrrrr}
    \toprule
    & \multicolumn{3}{c}{Per standpoint} & Walk \\
    \cmidrule(lr){2-4}
    Illumination & median & p90 & worst & median \\
    \midrule
    Isotropic          & 0.004 & 0.006 & 0.015 & 0.0042 \\
    Rooftop macro      & 0.024 & 0.040 & 0.066 & 0.0136 \\
    Street small cell  & 0.118 & 0.224 & 0.954 & 0.0343 \\
    \bottomrule
  \end{tabular}
\end{table}

The factor of about thirty between the isotropic and street rows is the reason
given in the body. The street law puts its whole mass into the lowest tens of
degrees of elevation, so most rays escape into directions the weight sets to
zero and contribute nothing.

Two caveats belong with \cref{tab:si-noise} and neither is small. It was
measured at six surface interactions on the 130\,m crop, which was the operating
point of the evidence ladder at the time, and not at the three-interaction
250\,m configuration of the headline run. The transfer to the headline
configuration is an argument rather than a measurement. And the numbers exist
only in the audit note that produced them. No replica file was written to
\texttt{outputs/}, and no per-standpoint value in any location file carries an
error bar of any kind.

Two smaller seed measurements are on disk and can be reproduced from the shipped
artefacts. The first sits inside the interaction budget sweep. At Korenmarkt on
the 250\,m crop, four standpoints were each traced eight times at 200\,000 rays
and three interactions, and the relative standard deviation of $\chi$ was
recorded. Its median over the four standpoints is 0.00081 isotropic, 0.00825
rooftop and 0.02177 street, which converts to 0.0035, 0.0358 and 0.0945\,dB. The
worst of the four is 0.00181, 0.00913 and 0.11778, or 0.0079, 0.0397 and
0.5115\,dB. The second sits inside the bystander study. Twelve standpoints on an
empty square were traced under three streams at 120\,000 rays and three
interactions, and the difference between two streams at the same standpoint has
a median absolute value of 0.0087, 0.0354 and 0.1123\,dB, a 95th percentile of
0.0240, 0.1194 and 0.6450\,dB and a maximum of 0.0274, 0.1399 and 0.7314\,dB.
The difference of two independent draws is $\sqrt{2}$ times a single one, so
these three sets of numbers agree with each other to the accuracy any of them
can claim, at three different ray counts and two different crops.

\subsection{Which comparisons are paired}

A comparison that changes only material or only the interaction budget keeps the
scene, the standpoints and the seeds, so the two sides share their launch set
and their difference is far quieter than either side. The evidence ladder, the
interaction budget sweep, the roulette on and off comparison, the two
illumination law comparisons and the material posterior ablation are all paired
in this sense. The illumination law comparisons are the strongest case, since
both laws are computed from one stored set of escapes and their difference
carries no sampling noise whatsoever.

Two comparisons are not paired and should not be read as though they were. The
crop sweep changes the mesh, so the two sides decorrelate at the first
interaction and the error on their difference is the independent one, about
$\sqrt{2}$ times the single-run figure. For the rooftop model that is roughly
0.019\,dB on a walk median, so the largest published crop correction of 4.80\,dB
is beyond argument while the smallest of 0.05\,dB is under three standard errors
and is not resolved at this ray count. The sub-street ablation changes the mesh
in the same way and inherits the same floor.

One number in the paper is reported as a seed average for this reason. The
street-cell shift across the evidence ladder ran from 0.079 to 0.269\,dB over
the eight streams, a factor of 3.4, and the single seed that had been quoted was
the lowest of the eight. The seed-averaged value is $0.167 \pm 0.022$\,dB, the
uncertainty being the standard error on a mean of eight. The isotropic and
rooftop shifts over the same eight streams are $0.347 \pm 0.001$ and
$0.383 \pm 0.006$\,dB, which are 105 and 22 standard errors from zero and were
not sensitive to the choice of seed.

\subsection{What the seed study does not cover}

Monte Carlo error is not the dominant uncertainty on a per-square number, and
adding rays does not touch the term that is. The standpoints themselves are a
sample of the square. Changing the ground datum by a fraction of a metre adds or
drops candidate positions, the ordering that selects 80 of them changes, and at
the nine sites whose datum was corrected only 4 to 24 of the 80 standpoints
survived the change. The per-square medians moved by 0.128, 0.061 and 0.188\,dB
rms and by up to 0.242, 0.116 and 0.326\,dB under the three illumination models.
Against a walk median Monte Carlo floor of a few hundredths of a decibel that is
more than an order of magnitude larger. Korenmarkt is the control, since its
standpoints did not change and its medians moved by less than the floor. The
consequence for reading the paper is stated in the body: a per-square median
carries roughly 0.13\,dB rms of standpoint sampling uncertainty, which does not
threaten a 3.71\,dB spread across eleven squares but does threaten any claim
that two adjacent squares differ.

\section{The eleven squares, one by one}
\label{sec:si-sites}

\cref{tab:si-sites} gives the numbers behind the headline figure of the paper,
per square. Every column comes from the same run, which is 11 sites at 80
standpoints, three surface interactions, the 250\,m crop and one common material
assignment, so what differs between rows is the shape of the square and nothing
else.

The standpoints are built the same way everywhere. A grid of 3\,m spacing is
laid over a disc of 90\,m radius about the centre of the square. A grid point is
kept when the surface under it is within 2.5\,m of the ground datum measured for
that square and faces up with a normal cosine of at least 0.85, when a head at
1.5\,m above it has at least 2\,m of clearance in every direction on a 96
direction probe, and when at least 1\% of the directions from that head reach
sky. The survivors are chained into a walk by nearest neighbour and 80 of them
are taken at even index spacing along that chain. The candidate count in
\cref{tab:si-sites} is the number that survived, and it runs from 801 at
Korenmarkt to 2114 at Krakow, so the 80 traced standpoints are a sample of
between 4 and 10\% of the walkable positions.

\begin{table*}[t]
  \centering
  \caption{The eleven squares at three surface interactions and the 250\,m
  crop, 80 standpoints each. $\chi$ columns are the median over the 80
  standpoints in dB. The spread column is the 95th percentile over the 5th
  within one square, also in dB.}
  \label{tab:si-sites}
  \footnotesize
  \begin{tabular}{lrrrrrrrrrr}
    \toprule
    & & Ground & Walk & \multicolumn{2}{c}{Isotropic} & \multicolumn{2}{c}{Rooftop macro} & \multicolumn{2}{c}{Street small cell} & Sky \\
    \cmidrule(lr){5-6}\cmidrule(lr){7-8}\cmidrule(lr){9-10}
    Square & Triangles & datum (m) & candidates & median & spread & median & spread & median & spread & fraction \\
    \midrule
    Brussels, Grand-Place       & 706\,720   & 65.57   & 1079 & $-6.53$ & 5.88 & $-10.06$ & 12.89 & $-21.77$ & 19.01 & 0.170 \\
    Ghent, Korenmarkt           & 617\,091   & 50.87   & 801  & $-5.35$ & 2.85 & $-7.89$  & 6.39  & $-18.45$ & 8.43  & 0.231 \\
    Krakow, Rynek Glowny        & 557\,526   & 250.94  & 2114 & $-4.16$ & 1.70 & $-5.72$  & 2.86  & $-16.35$ & 4.80  & 0.314 \\
    London, Trafalgar Square    & 783\,425   & 54.81   & 1619 & $-4.04$ & 4.04 & $-5.40$  & 4.55  & $-15.59$ & 8.31  & 0.330 \\
    Madrid, Plaza Mayor         & 720\,373   & 702.68  & 977  & $-4.64$ & 6.46 & $-8.40$  & 9.35  & $-21.01$ & 11.08 & 0.278 \\
    Mexico City, Zocalo         & 707\,812   & 2223.38 & 1343 & $-4.00$ & 2.19 & $-5.13$  & 5.06  & $-15.35$ & 9.16  & 0.324 \\
    Milan, Piazza del Duomo     & 779\,421   & 163.12  & 2026 & $-4.49$ & 2.80 & $-6.11$  & 2.55  & $-14.26$ & 6.20  & 0.286 \\
    New York, Times Square      & 1\,064\,389 & $-18.39$ & 1083 & $-7.71$ & 3.92 & $-8.96$  & 5.77  & $-12.19$ & 7.90  & 0.121 \\
    Prague, Old Town Square     & 664\,619   & 236.65  & 1495 & $-4.46$ & 4.29 & $-6.24$  & 4.95  & $-17.06$ & 7.27  & 0.288 \\
    Tokyo, Hachiko crossing     & 797\,615   & 51.60   & 1138 & $-6.32$ & 3.36 & $-10.02$ & 5.46  & $-17.92$ & 10.70 & 0.174 \\
    Toulouse, Place du Capitole & 1\,016\,440 & 191.23  & 1312 & $-5.28$ & 5.17 & $-6.54$  & 11.40 & $-17.59$ & 13.49 & 0.238 \\
    \bottomrule
  \end{tabular}
\end{table*}

Reading the table across gives the comparison the paper is built on. Between
squares the medians span 3.71\,dB isotropic, 4.93\,dB rooftop and 9.58\,dB
street. Within one square the spread runs 1.70 to 6.46\,dB isotropic, 2.55 to
12.89\,dB rooftop and 4.80 to 19.01\,dB street. Under isotropic illumination the
smallest within-square spread in the set, 1.70\,dB at Krakow, is already half
the whole between-square range, and the median within-square spread of 3.92\,dB
exceeds it. Under street illumination the ordering reverses, because there the
between-square range is 9.58\,dB against a median within-square spread of
8.43\,dB, and that is the case the paper reports as going the other way.

The between-square range is set by different pairs under the three models.
Isotropic runs from the Zocalo down to Times Square, rooftop from the Zocalo
down to Grand-Place, and street from Times Square down to Grand-Place, so Times
Square is at one end under two models and at the other end under the third. That
is the shape of the square doing the work. Its towers block the horizon, which
costs it under isotropic illumination, and they leave the high elevations
comparatively open, which pays under an illumination model whose weight sits low
and near.

The ground datum column is worth reading beside the rest, because it is a
measured surface height in each mesh's own vertical frame and not a free choice.
Its absolute value carries no meaning across squares and cancels in a ratio.
Krakow and Toulouse are the two squares whose datum was corrected during the
study, having previously sat 18.19\,m and 13.94\,m too high, which put their
whole walks on a roof. Coming down to the pavement costs Krakow 1.40\,dB
isotropic, 5.26\,dB rooftop and 12.97\,dB street, and Toulouse 1.29, 1.83 and
2.61\,dB. Every other square moved by less than 0.25\,dB isotropic under the
same fix, so the datum is not a detail at the two squares where it was wrong and
is not a worry anywhere else.

\section{Material assignment}
\label{sec:si-materials}

Every square in the headline run carries the same four-class assignment, so
material is held fixed while form varies. The class of a triangle follows from
its normal and its height. A triangle whose unit normal has $\hat{n}_z > 0.5$ is
ground when its centroid is within 4\,m of the ground datum and roof when it is
above that. A triangle with $\hat{n}_z < -0.5$ is a soffit. Everything else is a
facade. The datum is the one measured for that square, so the ground and roof
split is anchored to the pavement the walk stands on.

\cref{tab:si-materials} gives what each class is made of. The permittivity is
the ITU-R P.2040-4 power-law row for that material, $\epsilon_r = a f^b$ with
$f$ in GHz, evaluated at 15\,GHz, with the conductivity from the same row
carried into the imaginary part \cite{itu2040}. The roughness is an RMS height
that sets the specular and diffuse split.

The roughness column is the weakest part of the table and it is worth saying
why. P.2040-4 tabulates permittivity and conductivity and supplies no facade
roughness at all, so the roughness values come from a separate file that carries
no ITU standing. The functional form it uses, a coherent reflected fraction of
$\exp(-g^2)$ with $g = 4\pi s \cos\theta / \lambda$, is the one in ITU-R
P.2146-0, which is a sea-surface recommendation, so applying it to a brick wall
is an extrapolation. Of the four rows, the facade rests on a millimetre-wave
measurement of a brick wall and the soffit on laboratory metrology of cast
concrete, while the ground and the roof are extrapolated, and together those two
hold about half the triangle area of every square. The RMS height in
\cref{tab:si-materials} is also the effective one the tracer uses, which folds
any periodic step in the surface into the random part in quadrature. That
treatment overstates the coherent loss for a strictly periodic grating, so the
specular fractions in this study are a lower bound and the diffuse shares an
upper bound.

\begin{table}[t]
  \centering
  \caption{The four-class material assignment, common to all eleven squares.}
  \label{tab:si-materials}
  \footnotesize
  \setlength{\tabcolsep}{4pt}
  \begin{tabular}{llrl}
    \toprule
    Class & ITU-R row & RMS (mm) & Roughness evidence \\
    \midrule
    Ground  & asphalt concrete & 3.79 & extrapolated \\
    Facade  & brick            & 1.16 & mm-wave measurement \\
    Roof    & concrete         & 0.60 & extrapolated \\
    Soffit  & concrete         & 0.15 & laboratory metrology \\
    \bottomrule
  \end{tabular}
\end{table}

What that assignment amounts to in each square is set by how much area each
class holds, and that varies a great deal. \cref{tab:si-areas} gives the split.
Facade area runs from 43.9\% at the Zocalo to 72.6\% at Times Square, and ground
area runs from 10.8\% at Times Square to 25.7\% at Korenmarkt. The two are not
independent, because a square hemmed in by tall buildings has a large facade and
a small visible pavement. So although the material table is identical
everywhere, the material a ray meets is not, and part of what
\cref{tab:si-sites} shows as a difference in form is carried by that.

\begin{table}[t]
  \centering
  \caption{Share of triangle area per class at the 250\,m crop.}
  \label{tab:si-areas}
  \footnotesize
  \begin{tabular}{lrrrr}
    \toprule
    Square & Ground & Facade & Roof & Soffit \\
    \midrule
    Brussels, Grand-Place       & 0.197 & 0.501 & 0.251 & 0.051 \\
    Ghent, Korenmarkt           & 0.257 & 0.441 & 0.248 & 0.054 \\
    Krakow, Rynek Glowny        & 0.187 & 0.498 & 0.259 & 0.056 \\
    London, Trafalgar Square    & 0.152 & 0.531 & 0.212 & 0.104 \\
    Madrid, Plaza Mayor         & 0.149 & 0.529 & 0.281 & 0.041 \\
    Mexico City, Zocalo         & 0.201 & 0.439 & 0.246 & 0.114 \\
    Milan, Piazza del Duomo     & 0.187 & 0.577 & 0.147 & 0.089 \\
    New York, Times Square      & 0.108 & 0.726 & 0.095 & 0.072 \\
    Prague, Old Town Square     & 0.231 & 0.468 & 0.257 & 0.044 \\
    Tokyo, Hachiko crossing     & 0.144 & 0.608 & 0.163 & 0.086 \\
    Toulouse, Place du Capitole & 0.133 & 0.508 & 0.311 & 0.048 \\
    \bottomrule
  \end{tabular}
\end{table}

The image-derived route replaces the class of a triangle a panorama saw with one
read from the photograph, and leaves every other triangle on the geometric rule.
It reaches 3.2\% of triangle area at Korenmarkt at the 250\,m crop and less at
every other square, which is why the eleven-square comparison runs on the
geometric assignment and the image comparison is reported separately at one
square.

\section{Convergence sweeps}
\label{sec:si-convergence}

Two operating points had to be chosen, the crop radius and the number of surface
interactions, and both were measured rather than assumed. This section gives
both sweeps at more resolution than the body has room for.

\subsection{Crop radius}

\cref{tab:si-crop} sweeps the radius the scene is cut at while holding the
standpoints fixed inside the smallest crop, so every radius scores the same 32
observers and only the surroundings change. Every radius uses the same seed, so
the launch set is common to all nine and the comparison is paired at that level.

\begin{table}[t]
  \centering
  \caption{Crop radius sweep at Korenmarkt, 32 fixed observers, 200\,000 rays,
  three surface interactions. $\chi$ entries are means over the 32 observers.}
  \label{tab:si-crop}
  \footnotesize
  \setlength{\tabcolsep}{4pt}
  \begin{tabular}{rrrrrr}
    \toprule
    Radius (m) & Triangles & Isotropic & Rooftop & Street & Sky \\
    \midrule
    60  & 31\,706    & 0.30195 & 0.23900 & 0.14088 & 0.2351 \\
    100 & 97\,114    & 0.28731 & 0.17428 & 0.08365 & 0.2233 \\
    120 & 135\,640   & 0.28493 & 0.16344 & 0.06424 & 0.2216 \\
    130 & 157\,862   & 0.28755 & 0.16567 & 0.06253 & 0.2235 \\
    160 & 245\,112   & 0.28459 & 0.15131 & 0.02203 & 0.2215 \\
    200 & 390\,518   & 0.28376 & 0.14606 & 0.01376 & 0.2210 \\
    250 & 617\,091   & 0.28367 & 0.14559 & 0.01235 & 0.2209 \\
    300 & 903\,712   & 0.28365 & 0.14559 & 0.01218 & 0.2209 \\
    340 & 1\,159\,281 & 0.28364 & 0.14558 & 0.01209 & 0.2209 \\
    \bottomrule
  \end{tabular}
\end{table}

Convergence is read as the smallest radius from which every later step is under
half a decibel. On that rule the isotropic model and the sky fraction converge
at 60\,m, the rooftop model at 100\,m and the street model at 200\,m. The
ordering is the order in which the three models put their weight near the
horizon, which is also the order of how much scene each one needs, because a ray
leaving a standing head near the horizon travels a long way before it rises high
enough for a building of ordinary height to intercept it. The published radius
is 250\,m, past the last of the three.

The 130\,m row breaks the monotone trend in all three columns. The builds at 60,
100 and 120\,m are single precision and those from 130\,m up are double
precision, so the series carries a real discontinuity at 125\,m and the 130\,m
row should not be read as scene structure.

\subsection{Number of surface interactions}

\cref{tab:si-budget} sweeps the interaction budget at four squares against a
reference of eight interactions with roulette off, on the same standpoints and
the same seeds, so every row is a paired comparison and the differences carry
almost no Monte Carlo noise. The shipped operating point is three interactions
with roulette disabled.

\begin{table}[t]
  \centering
  \caption{Three surface interactions against a reference of eight, 40
  standpoints, 200\,000 rays, roulette off in both. Worst is the largest
  absolute change over the 40 standpoints, in dB. Truncated share is the power
  still in flight at the cut, as a share of the power that escaped.}
  \label{tab:si-budget}
  \footnotesize
  \setlength{\tabcolsep}{4pt}
  \begin{tabular}{lrrrrr}
    \toprule
    & \multicolumn{3}{c}{Worst change (dB)} & \multicolumn{2}{c}{Truncated share} \\
    \cmidrule(lr){2-4}\cmidrule(lr){5-6}
    Square & Isotr. & Rooftop & Street & median & worst \\
    \midrule
    Korenmarkt, 130\,m & 0.006 & 0.016 & 0.063 & 0.0037 & 0.0201 \\
    Korenmarkt, 250\,m & 0.006 & 0.016 & 0.049 & 0.0044 & 0.0201 \\
    Zocalo             & 0.006 & 0.031 & 0.071 & 0.0049 & 0.0267 \\
    Plaza Mayor        & 0.017 & 0.047 & 0.169 & 0.0048 & 0.0747 \\
    Grand-Place        & 0.017 & 0.071 & 0.384 & 0.0094 & 0.0798 \\
    \bottomrule
  \end{tabular}
\end{table}

No standpoint at any of the four squares moves by half a decibel when the budget
is cut from eight interactions to three. The worst single standpoint in the
whole set is 0.384\,dB under street illumination at Grand-Place, which is also
the most enclosed square in the study and the one with the largest truncated
share. The paper quotes 0.063\,dB as the worst change, and that is the
Korenmarkt figure at the 130\,m crop. It is not the worst across the four
squares that were measured, and \cref{tab:si-budget} is the fuller statement.

Truncation deletes power and never adds it, so a susceptibility at this budget
is a lower bound rather than a two-sided estimate. The share still in flight at
the cut is a few thousandths of the power that escaped at the median standpoint
and under a tenth at the worst.

One row of the sweep does cross half a decibel and it is not the shipped one.
With roulette switched on from the third interaction, two of the 40 Grand-Place
standpoints move past 0.5\,dB under street illumination, and they do so at the
four, six and three interaction budgets alike. Roulette is unbiased, so those
crossings are seed noise on the noisiest model at the most enclosed square
rather than a truncation effect, and they are absent from the roulette-off
configuration the study ships. The repository notes carry a claim that no
standpoint out of 40 moves half a decibel at any budget, and that is true at
Korenmarkt, where it was measured, and not at Grand-Place.

Russian roulette is disabled at the shipped budget and this was also measured
rather than preferred. At three interactions roulette can only kill a ray one
step before the hard cap drops it anyway, so there is no work to buy. Measured
over four Korenmarkt standpoints at the 250\,m crop with eight seeds each, the
relative standard deviation of $\chi$ is 0.00083, 0.00807 and 0.02218 with
roulette on against 0.00081, 0.00825 and 0.02177 with it off, which is the same
number to the accuracy four standpoints can resolve, and off is the faster of
the two at 101.9\,s against 106.4\,s for the same work. Roulette is unbiased
whichever way the switch is set, so it never was and is not now what bounds the
truncation.

\subsection{Where the power is}

Of the power launched from a standpoint at Korenmarkt, 79.0\% reaches a first
surface, 8.85\% reaches a second, 1.19\% reaches a third, and everything past
the third interaction adds up to 0.23\%. The remaining 21.0\% leaves the square
without touching anything, which is the sky fraction seen from the other side.

Those shares are pooled over the eight camera positions at the 130\,m crop
rather than over the pedestrian walk, which is the standpoint set every exposure
number in the paper uses. Pooled over the 40 walk standpoints instead, the
untouched share is 23.6\% at the 130\,m crop and 22.6\% at 250\,m. The cameras
stand in more open parts of the square than the average walk position, which is
the whole of the difference. The three shares should be read as the shape of the
budget and not as a per-standpoint figure.

A caution about a related number. One in the repository notes gives the
untouched share as 10.7\%, computed as one minus the sum of the per-depth
incident shares. Those shares are not a partition, because a ray that reaches
the third interaction was counted at the first and the second as well, so their
sum is not the touched fraction and one minus it is not the untouched one. The
correct figure is the first-depth complement given above.

\begin{thebibliography}{99}

\bibitem{itu2040}
International Telecommunication Union, ``Effects of building materials and
structures on radio-wave propagation in the range of 1 MHz to 450 GHz,''
Recommendation ITU-R P.2040-4, 2025.

\bibitem{numpy}
C.~R.~Harris \emph{et al.}, ``Array programming with NumPy,'' \emph{Nature},
vol. 585, no. 7825, pp. 357--362, 2020.


\bibitem{storn}
R.~Storn and K.~Price, ``Differential evolution: a simple and efficient
heuristic for global optimization over continuous spaces,'' \emph{J. Global
Optim.}, vol. 11, no. 4, pp. 341--359, 1997.

\bibitem{scipy}
P.~Virtanen \emph{et al.}, ``SciPy 1.0: fundamental algorithms for scientific
computing in Python,'' \emph{Nat. Methods}, vol. 17, no. 3, pp. 261--272, 2020.

\bibitem{mask2former}
B.~Cheng, I.~Misra, A.~G.~Schwing, A.~Kirillov, and R.~Girdhar,
``Masked-attention mask transformer for universal image segmentation,'' in
\emph{Proc. IEEE/CVF Conf. Comput. Vis. Pattern Recognit.}, 2022,
pp. 1290--1299.

\bibitem{vistas}
G.~Neuhold, T.~Ollmann, S.~Rota~Bul\`o, and P.~Kontschieder, ``The Mapillary
Vistas dataset for semantic understanding of street scenes,'' in \emph{Proc.
IEEE Int. Conf. Comput. Vis.}, 2017, pp. 4990--4999.

\bibitem{douglas}
D.~H.~Douglas and T.~K.~Peucker, ``Algorithms for the reduction of the number of
points required to represent a digitized line or its caricature,''
\emph{Cartographica}, vol. 10, no. 2, pp. 112--122, 1973.

\end{thebibliography}

\end{document}
